\vspace{-7pt}

\section{Conclusion}
\label{sec:conclusion}
In this paper, we showed that maximizing batch size alone is not sufficient for efficient LLM serving during the decode phase. Instead, we identified \textit{prefix homogeneity}, the extent of the prefix shared across \textit{all} requests in a batch, as a more important factor for improving memory locality and reducing redundant KV cache accesses. Higher batch sizes lead to better GPU utilization but also reduce the length of the prefix shared across all requests. Existing approaches ignore this aspect and also rely on expensive prefix detection mechanisms that introduce significant CPU overhead. To address these gaps, we proposed \Feather, a lightweight prefix-aware scheduler that balances this trade-off between batch size and prefix homogeneity. It uses a Chunked Hash Tree (CHT) for fast prefix detection and a reinforcement learning-based batching policy that decides when to stop batch formation based on the loss in prefix homogeneity due to the addition of a new request. Our results show that \Feather consistently improves end-to-end throughput across diverse workloads while keeping scheduling overhead low. As future work, we plan to extend this approach to distributed and multi-GPU settings. We are also looking into more adaptive and lightweight learning policies that can improve robustness across different workloads and deployment scenarios.